% TODO make hw stylesheet
\documentclass[10pt]{scrartcl}
\usepackage[a4paper, total={6in, 8in}]{geometry}
\usepackage[usenames,dvipsnames,svgnames]{xcolor}
\usepackage[shortlabels]{enumitem}
\usepackage[framemethod=TikZ]{mdframed}
\usepackage{amsmath,amssymb,amsthm}
\usepackage[colorlinks]{hyperref}
\usepackage{multicol}
\usepackage{tabularx}

\hypersetup{
	colorlinks=true,
	linkcolor=blue,
	filecolor=magenta,      
	urlcolor=magenta,
	pdftitle={MAT 362 HW}
}

\usepackage{microtype}
\usepackage{mathtools}
\usepackage[headsepline]{scrlayer-scrpage}
\usepackage{thmtools}
\usepackage[textsize=footnotesize]{todonotes}

\mdfdefinestyle{mdbluebox}{roundcorner=10pt,innerbottommargin=9pt,
	linecolor=blue,backgroundcolor=TealBlue!5,}
\declaretheoremstyle[headfont=\sffamily\bfseries\color{MidnightBlue},
mdframed={style=mdbluebox},]{thmbluebox}

\mdfdefinestyle{mdbloobox}{frametitlefont=\bfseries,innerbottommargin=8pt,
	nobreak=true,backgroundcolor=TealBlue!5,linecolor=blue,}
\declaretheoremstyle[headfont=\bfseries\color{MidnightBlue},
mdframed={style=mdbloobox},headpunct={\\[3pt]},postheadspace=0pt,]{thmbloobox}

\mdfdefinestyle{mdredbox}{frametitlefont=\bfseries,innerbottommargin=8pt,
	nobreak=true,backgroundcolor=Salmon!5,linecolor=RawSienna,}
\declaretheoremstyle[headfont=\bfseries\color{RawSienna},
mdframed={style=mdredbox},headpunct={\\[3pt]},postheadspace=0pt,]{thmredbox}

\mdfdefinestyle{mdgreenbox}{linecolor=ForestGreen,backgroundcolor=ForestGreen!5,
	linewidth=2pt,rightline=false,leftline=true,topline=false,bottomline=false,}
\declaretheoremstyle[headfont=\bfseries\sffamily\color{ForestGreen!70!black},
mdframed={style=mdgreenbox},headpunct={ --- },]{thmgreenbox}

\mdfdefinestyle{mdorangebox}{linecolor=RedOrange,backgroundcolor=RedOrange!5,
	linewidth=2pt,rightline=false,leftline=true,topline=false,bottomline=false,}
\declaretheoremstyle[headfont=\bfseries\sffamily\color{RedOrange!70!black},
mdframed={style=mdorangebox},headpunct={ --- },]{thmorangebox}

\mdfdefinestyle{mdblackbox}{linecolor=black,backgroundcolor=RedViolet!5!gray!5,
	linewidth=3pt,rightline=false,leftline=true,topline=false,bottomline=false,}
\declaretheoremstyle[mdframed={style=mdblackbox}]{thmblackbox}

\declaretheoremstyle[spaceabove=6pt,spacebelow=6pt]{basehead}

\declaretheorem[style=basehead,name=Problem]{problem}
\declaretheorem[style=basehead,name=Problem,numbered=no]{problem*}
\declaretheorem[style=thmbloobox,name=Setup]{setup} % for config geo
\declaretheorem[style=thmredbox,name=Example,sibling=problem]{example}
\declaretheorem[style=thmredbox,name=$\bigstar$ Problem,sibling=problem]{niceprob}
\declaretheorem[style=thmbluebox,name=Theorem,sibling=problem]{theorem}
\declaretheorem[style=thmbluebox,name=Corollary,sibling=theorem]{corollary}
\declaretheorem[style=thmbluebox,name=Proposition,sibling=theorem]{prop}
\declaretheorem[style=thmbluebox,name=Lemma,numberwithin=problem]{lemma}
\declaretheorem[style=thmbluebox,name=Theorem,numbered=no]{theorem*}
\declaretheorem[style=thmbluebox,name=Lemma,numbered=no]{lemma*}
\declaretheorem[style=thmgreenbox,name=Claim,numberwithin=problem]{claim}
\declaretheorem[style=thmgreenbox,name=Claim,numbered=no]{claim*}
\declaretheorem[style=thmblackbox,name=Remark,numberwithin=problem]{remark}
\declaretheorem[style=thmblackbox,name=Remark,numbered=no]{remark*}
\declaretheorem[style=thmgreenbox,name=Definition,sibling=theorem]{definition}
\declaretheorem[style=thmgreenbox,name=Definition,numbered=no]{definition*}
\declaretheorem[style=thmorangebox,name=Warning,sibling=theorem]{warning}
\declaretheorem[style=thmorangebox,name=Warning,numbered=no]{warning*}
\declaretheorem[style=thmorangebox,name=Note,numbered=no]{note}
\declaretheorem[style=thmblackbox,name=Exercise,sibling=problem]{exercise}
\declaretheorem[style=thmblackbox,name=Exercise,numbered=no]{exercise*}
\declaretheorem[style=thmblackbox,name=Question,sibling=problem]{question}
\declaretheorem[style=thmblackbox,name=Question,numbered=no]{question*}

\addtolength{\textheight}{3.14cm}
\providecommand{\ol}{\overline}
\providecommand{\ul}{\underline}
\providecommand{\eps}{\varepsilon}
\providecommand{\half}{\frac{1}{2}}
\providecommand{\dang}{\measuredangle} %% Directed angle
\providecommand{\CC}{\mathbb C}
\providecommand{\FF}{\mathbb F}
\providecommand{\NN}{\mathbb N}
\providecommand{\QQ}{\mathbb Q}
\providecommand{\RR}{\mathbb R}
\providecommand{\ZZ}{\mathbb Z}
\providecommand{\dg}{^\circ}
\providecommand{\ii}{\item}
\providecommand{\setdiff}{\smallsetminus}
\providecommand{\alert}[1]{{\sffamily\textbf{\textcolor{blue}{#1}}}}
\providecommand{\inv}{^{-1}}
\providecommand{\mb}[1]{\mathbf{#1}}
\providecommand{\dif}{\;d}
\newcommand*{\horzbar}{\rule[.5ex]{2.5ex}{0.5pt}}

\reversemarginpar

\newenvironment{soln}{\begin{proof}[Solution]}{\end{proof}}
\newenvironment{walkthrough}{\noindent \textbf{Walkthrough.}}{\medskip}
\newlist{walk}{enumerate}{3}
\setlist[walk]{label=\bfseries (\alph*)}

\providecommand{\pow}{\operatorname{Pow}}
\providecommand{\vocab}[1]{{\textbf{\textcolor{ForestGreen}{#1}}}}
\providecommand{\lcm}{\operatorname{lcm}}
\providecommand{\abs}[1]{\left|#1\right|}

\usepackage{asymptote}
\begin{asydef}
	defaultpen(fontsize(10pt));
	unitsize(1cm);
	size(8cm); // set a reasonable default
	usepackage("amsmath");
	usepackage("amssymb");
	settings.tex="pdflatex";
	settings.outformat="pdf";
	// Replacement for olympiad+cse5 which is not standard
	import geometry;
	// recalibrate fill and filldraw for conics
	void filldraw(picture pic = currentpicture, conic g, pen fillpen=defaultpen, pen drawpen=defaultpen)
	{ filldraw(pic, (path) g, fillpen, drawpen); }
	void fill(picture pic = currentpicture, conic g, pen p=defaultpen)
	{ filldraw(pic, (path) g, p); }
	// some geometry
	pair foot(pair P, pair A, pair B) { return foot(triangle(A,B,P).VC); }
	pair orthocenter(pair A, pair B, pair C) { return orthocentercenter(A,B,C); }
	pair centroid(pair A, pair B, pair C) { return (A+B+C)/3; }
	// cse5 abbreviations
	path CP(pair P, pair A) { return circle(P, abs(A-P)); }
	path CR(pair P, real r) { return circle(P, r); }
	pair IP(path p, path q) { return intersectionpoints(p,q)[0]; }
	pair OP(path p, path q) { return intersectionpoints(p,q)[1]; }
	path Line(pair A, pair B, real a=0.6, real b=a) { return (a*(A-B)+A)--(b*(B-A)+B); }
	// cse5 more useful functions
	picture CC() {
		picture p=rotate(0)*currentpicture;
		currentpicture.erase();
		return p;
	}
	pair MP(Label s, pair A, pair B = plain.S, pen p = defaultpen) {
		Label L = s;
		L.s = "$"+s.s+"$";
		label(L, A, B, p);
		return A;
	}
	pair Drawing(Label s = "", pair A, pair B = plain.S, pen p = defaultpen) {
		dot(MP(s, A, B, p), p);
		return A;
	}
	path Drawing(path g, pen p = defaultpen, arrowbar ar = None) {
		draw(g, p, ar);
		return g;
	}
\end{asydef}

\usepackage{mathrsfs}

\ihead{\footnotesize MAT 362 HW04}
\ohead{\footnotesize Last compiled \today}

\title{\vspace{-2em}MAT 362 HW04}
\author{\large Aditya Pahuja}
\date{\large Due 03/07}
\begin{document}
\maketitle

\begin{problem*}2.3.6
	Prove that the definition of a differentiable map between surfaces
	does not depend on the parametrization chosen.
\end{problem*}

\begin{soln}
	Let $f\colon S_1\to S_2$ be differentiable at $p\in S_1$
	with respect to parametrizations
	$\mb x_1\colon U_1\subseteq\RR^2\to S_1$ and
	$\mb x_2\colon U_2\subseteq\RR^2\to S_2$
	(where $p\in \mb x_1(U_1)$, $f(p)\in \mb x_2(U_2)$); that is,
	\[\mb x_2\inv \circ f\circ\mb x_1\colon U_1\to U_2\]
	is differentiable at $\mb x_1\inv(p)$.
	
	Then, given parametrizations $\ol{\mb x}_1\colon V_1\subseteq\RR^2\to S_1$
	and $\ol{\mb x}_2\colon V_2\subseteq\RR^2\to S_2$
	such that $p\in\ol{\mb x}_1(V_1)$ and $f(p)\in\ol{\mb x}_2(V_2)$,
	we want to show that
	\[\ol{\mb x}_2\inv\circ f\circ\ol{\mb x}_1\colon V_1\to V_2\]
	is differentiable at $\ol{\mb x}_1\inv(p)$.
	This follows from the ``change of parameters'' maps
	\[\ol{\mb x}_1\inv\circ\mb x_1\colon \mb x_1\inv(W_1)\to
	\ol{\mb x}_1\inv (W_1),\quad
	\ol{\mb x}_2\inv\circ\mb x_2\colon \mb x_2\inv(W_2)\to
	\ol{\mb x}_2\inv (W_2)\]
	being diffeomorphisms
	(where $W_1$ is $\mb x_1(U_1)\cap\ol{\mb x}_1(V_1)$ and
	$W_2$ is $\mb x_2(U_2)\cap\ol{\mb x}_2(V_2)$), as we can write
	\[\ol{\mb x}_2\inv\circ f\circ\ol{\mb x}_1
	= (\ol{\mb x}_2\inv\circ\mb x_2)\circ
	(\mb x_2\inv \circ f\circ\mb x_1)\circ
	(\mb x_1\inv\circ\ol{\mb x}_1\inv)\]
	as a composition of differentiable functions.
\end{soln}

\begin{problem*}2.3.7.
	Prove that the relation ``$S_1$ is diffeomorphic to $S_2$''
	is an equivalence relation on the set of regular surfaces.
\end{problem*}

\begin{soln}
	The relation is reflexive because the identity map
	is trivially a diffeomorphism $S\to S$ for all regular surfaces $S$.
	
	The relation is symmetric by definition:
	if there exists a diffeomorphism $f\colon S_1\to S_2$
	that is differentiable and has differentiable inverse,
	then $f\inv\colon S_2\to S_1$ is
	differentiable with differentiable inverse.
	
	The relation is transitive because
	given surfaces $S_1$, $S_2$, $S_3$ with diffeomorphisms
	$f\colon S_1\to S_2$ and $g\colon S_2\to S_3$,
	the composition $g\circ f\colon S_1\to S_3$ is also a diffeomorphism:
	$g\circ f$ and its inverse $f\inv\circ g\inv$ are both differentiable
	because compositions of differentiable functions are differentiable.

	Having checked all three conditions for a relation to be
	an equivalence relation, we can conclude that being diffeomorphic
	is an equivalence relation on the set of regular surfaces.
\end{soln}

\begin{problem*}2.3.11.
	Prove that the rotations of a surface of revolution $S$ about its axis
	are diffeomorphisms of $S$.
\end{problem*}

\begin{soln}
	Suppose WLOG that the $z$-axis is the rotation axis of $S$.
	Then counterclockwise rotation about this axis by some fixed $\theta$
	is the linear map given by the matrix
	\[R_\theta = \begin{pmatrix}
		\cos\theta & -\sin\theta & 0\\
		\sin\theta & \cos\theta & 0\\
		0 & 0 & 1
	\end{pmatrix}\]
	in the standard basis.
	Linear maps are differentiable, and
	since $R_\theta$ is invertible, it has an inverse
	which is of course also differentiable.
	Thus rotations about an axis are diffeomorphisms of $\RR^3$.
	Then, since $S$ is a surface of revolution whose axis of rotation
	is the $z$-axis, the restriction $R_\theta|_S$ is a bijection on $S$
	and retains its differentiability properties (cf. Example 3 in the book),
	so rotations $R_\theta|_S$ are a diffeomorphisms of $S$.
\end{soln}

\begin{problem*}2.3.15.
	Let $C$ be a regular curve and let $\alpha\colon I\subseteq\RR\to C$,
	$\beta\colon J\subseteq\RR\to C$ be two parametrizations of $C$
	in a neighborhood of $p\in\alpha(I)\cap \beta(J) = W$. Let
	\[h = \alpha\inv\circ\beta\colon\beta\inv(W)\to\alpha\inv(W)\]
	be the change of parameters. Prove that
	\begin{enumerate}[(a)]
		\ii $h$ is a diffeomorphism.
		\ii The absolute value of the arc length of $C$ in $W$
		does not depend on which parametrization is chosen to define it;
		that is,
		\[\abs{\int_{t_0}^t |\alpha'(t)|\dif t} =
		\abs{\int_{\tau_0}^\tau |\beta'(\tau)|\dif \tau},\quad
		t=h(\tau),\;t\in I,\;\tau\in J\]
	\end{enumerate}
\end{problem*}

\begin{soln}
	(a) First, $h$ must be a homeomorphism because it is a composition
	of homeomorphisms $\alpha\inv$ and $\beta$.
	Let $r\in\beta\inv(W)$ and set $q = h(r)$.
	Since $\alpha(t) = (x(t), y(t), z(t))$ is a parametrization,
	we can assume, by renaming the axes if necessary, that $x'(q)\ne 0$.
	
	We extend $\alpha$ to a map $F\colon I\times\RR^2\to\RR^3$ defined by
	\[F(t, u, v) = (x(t), y(t) + u, z(t) + v).\]
	It is clear that $F$ is differentiable and that the restriction
	$F|_{I\times\{(0,0)\}} = \alpha$. Moreover,
	\[\det(dF_q) = \begin{vmatrix}
		x'(q) & 0 & 0\\
		y'(q) & 1 & 0\\
		z'(q) & 0 & 1
	\end{vmatrix} = x'(q)\ne 0.\]
	Thus, we can apply the inverse function theorem, which guarantees
	the existence of a neighborhood $M$ of $\alpha(q)$ in $\RR^3$
	such that $F\inv$ exists and is differentiable in $M$.
	
	By the continuity of $\beta$, there exists a neighborhood $N$ of $r$ in $J$
	such that $\beta(N)\subseteq M$. When restricted to $N$,
	$h|_N = F\inv\circ \beta|_N$ is a composition of differentiable maps,
	so it is differentiable at $r$. Since $r$ is arbitrary,
	$h$ is differentiable on $\beta\inv(W)$.
	
	The same argument applies to show $h\inv$ is differentiable,
	so $h$ is a diffeomorphism.
	
	(b) Writing $\beta = \alpha\circ h$, we have
	\[\int_{\tau_0}^\tau |\beta'(\tau)|\dif\tau =
	\int_{\tau_0}^\tau |(\alpha\circ h)'(\tau)|\dif\tau =
	\int_{\tau_0}^\tau |\alpha' (h(\tau))|\cdot |h'(\tau)|\dif\tau.\]
	Since $h$ is differentiable with differentiable inverse,
	$h$ is continuous and $h'$ can never be zero,
	meaning $h$ is either always positive or always negative.
	This lets us get rid of the absolute value around $h'(t)$ to compute
	\[\abs{\int_{\tau_0}^\tau |\alpha' (h(\tau))|\cdot |h'(\tau)|\dif\tau}
	= \abs{\int_{\tau_0}^\tau |\alpha' (h(\tau))|\cdot h'(\tau)\dif\tau}
	= \abs{\int_{t_0}^t |\alpha' (t)| \dif t}\]
	by $u$-substitution.
\end{soln}

\begin{problem*}2.4.3.
	Show that the equation of the tangent plane of a surface
	which is the graph of a differentiable function $z = f(x,y)$,
	at the point $p_0 = (x_0,y_0)$, is given by
	\[z = f(x_0,y_0) + f_x(x_0,y_0)(x - x_0) + f_y(x_0,y_0)(y - y_0).\]
\end{problem*}

\begin{soln}
	The surface admits a parametrization $\mb x\colon\RR^2\to S$ given by
	\[\mb x(x, y) = (x, y, f(x,y)).\]
	As per Proposition 1, the set of tangent vectors to $S$ at
	$(x_0,y_0,f(x_0,y_0))$
	is the image of the derivative
	\[d\mb x_{p_0} = \begin{pmatrix}
		1&0\\
		0&1\\
		f_x(x_0,y_0)&f_y(x_0,y_0).
	\end{pmatrix}\]
	In particular a normal vector to the tangent plane is
	\[\begin{pmatrix}
		1\\0\\f_x(x_0,y_0)
	\end{pmatrix}\wedge\begin{pmatrix}
		0\\1\\f_y(x_0,y_0)
	\end{pmatrix} =\begin{pmatrix}
		-f_x(x_0,y_0)\\-f_y(x_0,y_0)\\1
	\end{pmatrix}\]
	which tells us that the tangent plane passing through
	$(x_0,y_0,f(x_0,y_0))$ has equation
	\[z-f(x_0,y_0) - f_x(x_0,y_0)(x-x_0) - f_y(x_0,y_0)(y - y_0) = 0,\]
	which rearranges to the given equation.
\end{soln}

\begin{problem*}2.4.8.
	Prove that if $L\colon\RR^3\to\RR^3$ is a linear map and $S\subseteq\RR^3$
	is a regular surface invariant under $L$, i.e. $L(S)\subseteq S$,
	then the restriction $L|_S$ is a differentiable map and
	\[dL_p(w) = L(w)\]
	for $p\in S$, $w\in T_p(S)$.
\end{problem*}

\begin{soln}
	It is obvious that $L|_S$ is a differentiable map
	because $L$, being a linear map, is differentiable
	as a function $\RR^3\to\RR^3$ and $dL_p = L$ for all $p\in\RR^3$
	(again, as a function $\RR^3\to\RR^3$).
	For a given $w\in T_p(S)$, let $\alpha\colon(-\eps,\eps)\to S$
	be a differentiable parametrized curve with $\alpha(0) = p$,
	$\alpha'(0) = w$. Then, by the chain rule,
	\[dL_p(w) = (L\circ\alpha)'(0) = dL_{\alpha(0)}\cdot \alpha'(0)
	= L(w)\]
	where the multiplication $dL_{\alpha(0)}\cdot\alpha'(0)$
	is matrix multiplication.
\end{soln}

\begin{problem*}2.4.15.
	Show that if all normals to a connected surface pass through a fixed point,
	the surface is contained in a sphere.
\end{problem*}

\begin{soln}
	Let $S$ be the surface and suppose WLOG (say after a translation)
	that the fixed point is the origin.
	
	\begin{claim*}
		For any $p\in S$, if $V\subseteq S$ is
		the image of an open ball in $\RR^2$ containing $\mb x\inv(p)$,
		then $\|p\| = \|x\|$ for every $x\in V$.
	\end{claim*}
	\begin{proof}
		Let $U = \mb x\inv(V)$ be the open ball mentioned in the claim.
		For any $x\in V$, let $\alpha\colon I =(-\eps,1+\eps)\to \RR^2$
		be the function $\alpha(t) = (1-t)q + tw$ for 
		$q := \mb x\inv(p)$ and $w := \mb x\inv(x)$,
		where we choose $\eps$ small enough for $\alpha(I)\subseteq U$ to hold
		(such an $\eps$ exists because $U$ we can find
		open balls around $q$ and $w$ that are contained in $U$).
		Notably, we are making good use of the convexity of $U$.
		
		Then, $\alpha$ determines a regular curve through $p$ and $x$
		on the surface $S$ via the composed map $I\to U\to S$ defined by
		\[t\stackrel{\alpha}{\mapsto} (1-t)q + tw\stackrel{\mb x}{\mapsto}
		\mb x((1-t)q + tw),\]
		which we can easily verify is a differentiable homeomorphism
		whose derivative is never zero (really we just need to check this
		for $\alpha$, but that's trivial).
		
		Let $\beta = \mb x\circ\alpha$ be this composed map.
		On the interval $I$, we then have
		\[\beta'(t)\cdot\beta(t) = 0\]
		for all $t$ since $\beta(t)$ is by assumption
		normal to $S$ and $\beta'(t)$ is a tangent vector;
		but this is the derivative of $\beta(t)\cdot\beta(t) =
		\|\beta(t)\|^2$, so $\|\beta(t)\|^2$ must be constant
		as $t$ varies across $I$, hence $\|p\| = \|\beta(0)\| = \|\beta(1)\|
		= \|x\|$ for all $x\in V$.
	\end{proof}
	
	\pagebreak
	\begin{claim*}
		For any two points $a$, $b$ in a connected open subset $U$
		of an arbitrary metric space,
		there is a finite sequence of open balls $B_0$, $B_1$, \dots, $B_k$
		such that $B_i\subseteq U$ for all $i$,
		$a\in B_1$, $b\in B_k$, and $B_{j-1}\cap B_j\ne\varnothing$
		for all $j=1,\dots,k$.
	\end{claim*}

	\begin{proof}
		For a given point $a\in U$, we say that $b\in U$ is
		\emph{$a$-reachable} if such a sequence of open balls
		(which we will refer to as a ``chain from $a$ to $b$'')
		exists for $a$ and $b$. We will show that $R = U$, i.e.
		every point of $U$ is $a$-reachable.
		
		For any $x\in R$ we have a chain from $a$ to $x$,
		in which the terminal ball, say $B(y, r)$, contains $x$.
		Then, we can find an open ball centered at $x$, such as
		$B(x, r - B(y, r))$, contained in $B(y, r)$ ---
		that is, for any $x\in R$, there is an open ball centered at $x$
		that is contained in $R$. Thus, $R$ is open.
		
		On the other hand, $R$ is closed with respect to $U$.
		Indeed, if $x\in U$ is a limit point of $R$,
		then there is an open ball $B$ centered at $x$ that intersects $U$
		at some point $y\ne x$. Then we have a chain
		$B_0$, \dots, $B_n$ from $a$ to $y$, so $B_0$, \dots, $B_n$, $B$
		is a chain from $a$ to $x$, hence $x\in R$.
		Thus, $R$ is closed in $U$ because it contains its limit points.
		
		Since path connectedness implies open set connectedness
		(there are no nontrivial clopen sets in $U$),
		we can conclude that, because $R$ is both closed and open,
		it is either empty or equal to $U$, but clearly $R$
		is nonempty because $a\in R$, hence $R = U$ as desired.
	\end{proof}
	
	To finish, let $a$ and $b$ be two arbitrary points in $S$.
	Suppose for a moment that $a$ and $b$ are not negative scalar multiples
	of one another, i.e. there is no positive real $\lambda$
	such that $a = -\lambda b$. Then we can apply a rotation,
	sending $S$ to $S'$, so that $a$ and $b$ both have positive $z$-coordinate.
	
	Let $S'_+$ be the intersection of $S'$ with $\{(x,y,z) : z>0\}$.
	The projection $\pi|_{S'_+}$ of $S'_+$ onto the $xy$-plane
	must be injective, else there exists a point at which
	the normal vector does not pass through the origin because
	the tangent plane will somewhere be perpendicular to the $xy$-plane
	(at some point where the distance to the origin is locally maximized);
	this can only occur at points where $S'$ intersects the $xy$-plane.
	
	Thus undoing the projection is a parametrization
	$\pi\inv\colon \pi(S'_+)\to S'_+$. Now $\pi(a)$ and $\pi(b)$
	can be connected by a chain of open balls contained in $\pi(S'_+)$,
	at which point applying the first claim to each open ball in the chain
	and using the fact that consecutive balls in the sequence intersect
	shows that $\|a\| = \|b\|$.
	
	The only other consideration is when $a = -\lambda b$ for a
	positive real $\lambda$, whence the rotation to obtain $S'$ doesn't exist.
	In this case we pick an arbitrary third point $c$ on the surface
	that is not collinear with $a$ and $b$ (just pick $c$
	to be in some neighborhood $V\subseteq S$ of $a$)
	and then apply the logic of the previous paragraph to get
	$\|a\| = \|c\| = \|b\|$.
\end{soln}

\begin{problem*}2.4.18.
	Prove that if a regular surface $S$ meets a plane $P$ in a single point $p$,
	then this plane coincides with the tangent plane of $S$ at $p$.
\end{problem*}

\begin{soln}
	Apply a rigid motion so that $p$ maps to the origin
	and $P$ gets sent to the $xy$-plane.
	Then, on a small enough neighborhood $U\subseteq S$ of $p$,
	the $z$-coordinates of all $x\in U$ have the same sign, say WLOG positive
	(except for a $z$-coordinate of zero at $p$).
	For any differentiable parametrized curve $\alpha\colon(-\eps,\eps)\to S$
	on $S$ with $\alpha(0) = p$, the $z$ component function thus
	attains a local minimum at $0$, meaning $\alpha'(0)$ has a $z$-coordinate
	of zero, i.e. it lies in the $xy$-plane.
	This holds for every such $\alpha$, so every tangent vector at $p$
	lies in $P$, meaning $P$ is the tangent plane of $S$ at $p$.
\end{soln}



























\end{document}